% ============================================================
% Capitolo 5 — PROGRESS.md e il cambio di fase
% ============================================================

\chapter{\texttt{PROGRESS.md} e il cambio di fase}
\label{ch:progress-md}

\texttt{\_CONTEXT.md} fotografa il presente: dove sei adesso, cosa stai facendo, quali vincoli rispettare. Ma non risponde a domande come ``perché abbiamo scelto Prisma invece di Drizzle?'' o ``cosa è successo nella sessione di giovedì scorso quando il deploy è andato storto?''

Quelle risposte vivono in \texttt{PROGRESS.md}.

\section{La differenza tra context e progress}
\label{sec:05-differenza}

\begin{table}[htbp]
  \centering
  \begin{tabular}{ll}
    \toprule
    \textbf{\texttt{\_CONTEXT.md}} & \textbf{\texttt{PROGRESS.md}} \\
    \midrule
    Stato attuale & Storia del progetto \\
    Si sovrascrive & Cresce solo in avanti \\
    Letto sempre dall'agente & Letto su richiesta \\
    Fotografa dove sei & Registra da dove sei venuto \\
    Conciso per design & Può essere lungo \\
    \bottomrule
  \end{tabular}
  \caption{Differenza tra \texttt{\_CONTEXT.md} e \texttt{PROGRESS.md}}
  \label{tab:05-context-vs-progress}
\end{table}

Pensa a \texttt{\_CONTEXT.md} come al cruscotto di un'auto --- mostra la velocità attuale, il carburante, la temperatura. Pensa a \texttt{PROGRESS.md} come al libretto di manutenzione --- registra ogni intervento, ogni anomalia, ogni decisione presa nel tempo.

\section{Struttura di \texttt{PROGRESS.md}}
\label{sec:05-struttura}

Il file cresce dall'alto verso il basso: le voci più recenti sono in cima.

\begin{lstlisting}[language=,caption={Esempio di PROGRESS.md per TaskFlow API}]
# PROGRESS — TaskFlow API

---

## 2026-07-03 — T-009 Auth middleware refactor

- Estratto il middleware JWT in src/middleware/auth.ts
- Identificato: Zod eseguiva dopo auth — corretto l'ordine

DECISIONE: middleware condiviso invece di per-route.
Motivazione: riduce duplicazione, comportamento coerente.

---

## 2026-06-27 — T-007 POST /tasks (70%)

- POST /tasks: Zod, auth JWT, RFC 7807, 4 test passanti
- Migration: add-task-userid-index (dev only)
- HALT trigger su schema.prisma → confermato e proceduto
\end{lstlisting}

\section{Cosa scrivere in ogni voce}
\label{sec:05-cosa-scrivere}

Una voce utile contiene tre tipi di informazioni:

\textbf{Fatto} --- cosa è stato prodotto o modificato concretamente.

\textbf{Scoperta} --- qualcosa di non ovvio trovato durante il lavoro: bug latenti, comportamenti inaspettati delle librerie, assunzioni sbagliate.

\textbf{Decisione} --- una scelta architetturale, di design o di processo, con la motivazione.

Le scoperte e le decisioni sono le più preziose. Sono le cose che non si deducono dal codice --- sono il contesto dietro al codice.

\section{Come \texttt{PROGRESS.md} accelera le sessioni future}
\label{sec:05-accelerazione}

Giulia --- non presente giovedì durante il refactor --- apre una sessione e chiede:

\begin{quote}
``Carica \texttt{\_CONTEXT.md} e \texttt{PROGRESS.md}, poi dimmi cosa è cambiato nell'autenticazione nelle ultime due settimane.''
\end{quote}

L'agente risponde in trenta secondi con un riepilogo completo del refactor, la decisione architetturale presa e il lavoro ancora da fare. Giulia ha il contesto senza aver bisogno di leggere la git history o chiedere a Lorenzo.

\section{Il cambio di fase}
\label{sec:05-cambio-fase}

Una delle voci più importanti è il cambio di fase:

\begin{lstlisting}[language=,caption={Voce di cambio fase in PROGRESS.md}]
## 2026-07-15 — CAMBIO FASE: Design → Implementation

Stack confermato: Node.js 22 / Fastify 5 / PostgreSQL 16
Decisioni architetturali prese:
- Paginazione: cursor-based (AD-01 chiusa)
- Cache: nessuna al primo rilascio (AD-02 rimandata)

API contract: 12 endpoint (vedi docs/api-contract.md)
Schema DB v1 migrato in dev

Primo task Phase 3: T-001 Setup struttura e dipendenze
\end{lstlisting}

Questa voce serve da spartiacque: tutto sopra è storia di design, tutto sotto è storia di implementazione.

\section{Aggiornare \texttt{PROGRESS.md} senza sforzo}
\label{sec:05-aggiornamento}

Alla fine di ogni sessione usa \texttt{@checkpoint}:

\begin{quote}
\texttt{@checkpoint} --- proponi una voce per PROGRESS.md
\end{quote}

L'agente produce una voce pronta. Tu la revisioni, aggiungi il contesto che solo tu conosci, e la incolli nel file. Trenta secondi di lavoro.

\section{\texttt{PROGRESS.md} in team}
\label{sec:05-team}

In un team di tre persone, \texttt{PROGRESS.md} diventa il canale di comunicazione asincrona sul progetto. Ogni sviluppatore scrive la propria voce alla fine della sessione.

\begin{warnbox}
\texttt{PROGRESS.md} non sostituisce la comunicazione di team. Non è una board kanban, non è un sistema di ticket. È un diario tecnico che integra il contesto degli agenti AI. Tenete separati i due strumenti.
\end{warnbox}

\section*{\LabelSummary}

\begin{itemize}
  \item \texttt{PROGRESS.md} registra la storia del progetto: fatto, scoperte, decisioni con motivazione.
  \item Si aggiorna a fine sessione con \texttt{@checkpoint}; tu revisioni e incolli.
  \item Le \textbf{scoperte} e le \textbf{decisioni} (con motivazione) sono le voci più preziose.
  \item Il cambio di fase è una voce speciale che documenta lo spartiacque tra una fase e l'altra.
  \item In team, permette agli agenti di ricostruire il contesto delle sessioni di altri membri.
\end{itemize}
